% !TeX program = lualatex
% ================================================================
% Polish Tier 3 Vocabulary Version of BearingsChallenge
%
% Generated by the server using the following settings:
%
%   language            Polish (pol)
%   text direction      left to right
%   font                Noto Serif
%   fallback font       Noto Serif
%   built from          latex/worksheets/geometry/bearings/bearingsChallenge.tex
%   vocabulary key      latex/worksheets/geometry/bearings/bearingsChallenge/L2/pol/bearingsChallenge_polishKey.tex
%   tier 3 vocabulary   green   RGB 0 128 0
%   English text        black   RGB 0 0 0
%   Polish text         purple  RGB 112 48 160
%   layout macros       version 3
%
% Please don't edit any information in this comment block - the server
% compares it against the current settings and rebuilds this file when
% they differ, so changes here would be overwritten. However, feel free
% to make updates to the LaTeX below if you want to edit it.
% ================================================================

\documentclass{article}
% ================================================================
% Copied from the English sheet this was translated from, so that
% anything it sets up for itself still works here
% ================================================================
\usepackage{gensymb}
\usepackage{amsfonts}
\usepackage{graphicx} % Required for inserting images
\usepackage{tabularx}

\title{Bearings Challenge}

\date{}

% Characters this language's font has not got are borrowed from another,
% rather than being silently dropped - see docs/translated-sheet-latex.md
\directlua{%
  if luaotfload and luaotfload.add_fallback then
    luaotfload.add_fallback("eallatinfallback",
      { "Noto Serif:mode=harf;" })
  else
    tex.error("this luaotfload is too old to register a fallback font")
  end
}
\usepackage[english]{babel}
\babelprovide[import]{polish}
\babelfont[polish]{rm}[Renderer=Harfbuzz, BoldFont={Noto Serif}, ItalicFont={Noto Serif}, BoldItalicFont={Noto Serif}, RawFeature={fallback=eallatinfallback}]{Noto Serif}
\babelfont[polish]{sf}[Renderer=Harfbuzz, BoldFont={Noto Serif}, ItalicFont={Noto Serif}, BoldItalicFont={Noto Serif}, RawFeature={fallback=eallatinfallback}]{Noto Serif}
\newcommand{\ealtext}[1]{\foreignlanguage{polish}{#1}}
\newcommand{\ealtextblock}[1]{{\raggedright\foreignlanguage{polish}{#1}\par}}
% ================================================================
% EAL layout helpers
% ================================================================
\usepackage{xcolor}

\definecolor{ealkeycolour}{RGB}{0,128,0}
\definecolor{ealenglish}{RGB}{0,0,0}
\definecolor{ealtwocolour}{RGB}{112,48,160}

\newsavebox{\ealboxen}
\newsavebox{\ealboxtr}
\newlength{\ealglosswd}
\newlength{\ealglossraise}
\setlength{\ealglossraise}{1.15em}

% a tier 3 word, in English and in translation alike
\newcommand{\ealkey}[1]{\textcolor{ealkeycolour}{#1}}
\newcommand{\ealkeytr}[1]{\textcolor{ealkeycolour}{\ealtext{#1}}}

% One translated word above one English word. Built from kernel
% primitives, never a tabular, which would break wherever the sheet
% loads array, booktabs or tabularx. The box is as wide as the wider of
% the two, so a long translation pushes its neighbours aside instead of
% overlapping them, and it claims its full height so a table row or a
% TikZ node makes room for it.
\newcommand{\ealgloss}[2]{%
  \sbox{\ealboxen}{\ealkey{#1}}%
  \sbox{\ealboxtr}{\small\textcolor{ealtwocolour}{\ealtext{#2}}}%
  \setlength{\ealglosswd}{\wd\ealboxen}%
  \ifdim\wd\ealboxtr>\ealglosswd\setlength{\ealglosswd}{\wd\ealboxtr}\fi
  \leavevmode
  \makebox[\ealglosswd][c]{%
    \makebox[0pt][c]{\usebox{\ealboxen}}%
    \makebox[0pt][c]{\raisebox{\ealglossraise}{\usebox{\ealboxtr}}}%
  }%
}

% opens up line spacing around a block of glosses, so the raised
% translations sit clear of the line above
\newenvironment{ealglossed}{\par\linespread{2.1}\selectfont\ignorespaces}{\par}

% a whole translated sentence above its English counterpart. block level,
% so it does not belong inside a tabular cell
\newcommand{\ealpara}[2]{%
  \par\addvspace{0.5\baselineskip}%
  {\color{ealtwocolour}\ealtextblock{#1}}%
  \penalty200
  {\color{ealenglish}#2\par}%
  \addvspace{0.5\baselineskip}%
}

\begin{document}
\maketitle

Bob knows the location of some buried gold, it lies exactly $100$km away, on a line exactly $123.5 \degree$ clockwise from North.

\begin{ealglossed}
Bob has an electronic compass that can direct him on \ealgloss{bearings}{azymuty}, with no error, unfortunately, it will only accept 3 digit \ealgloss{bearings}{azymuty}, with no \ealgloss{fractions}{ułamki} of \ealgloss{degrees}{stopnie}.
\end{ealglossed}

To find the gold, Bob must arrive within $1$m of its location.

\section{Introduction}

Assume all travellers can perfectly measure their distance travelled.

\begin{ealglossed}
\begin{enumerate}
    \item Using a \ealgloss{scale}{skalować} of $1$cm $:$ $20$km, sketch a map for Bob, label the \ealgloss{angles}{kąty} and distances
    \item Determine how Bob can locate the gold by travelling on a \ealgloss{bearing}{azymut} of $180\degree$ followed by $090\degree$
    \item Bob can travel $20$km a day, how long will it take him to travel to the gold?
    \item Rob has stumbled upon the map to the gold, and intends to plot a course to arrive before Bob, travelling Southeast, then East. Armed with the same model of electronic compass, and travelling at the same speed, how much earlier will he arrive?
    \item You want to beat both Bob and Rob to the treasure, but can also only follow $3$ digit \ealgloss{bearings}{azymuty}. And now they both have a $7$ hour head start! Covering $20$km per day can you find a way to arrive at the gold first? is your route the quickest possible?
\end{enumerate}
\end{ealglossed}

\newpage

\section{Challenge I}

Previously we assumed that the travellers could measure their distances with unlimited precision - in reality that is not the case. Luckily, you, Bob and Rob have downloaded an app that will ping every time you travel exactly $1$km in a straight line.

Therefore - all sections of journeys must be exact multiples of $1$km!

\begin{ealglossed}
\begin{enumerate}
    \item Determine how far off the gold's location you, Bob and Rob will be if you \ealgloss{round}{zaokrąglisz} each leg of your planned journeys to the nearest kilometre. Thus show that the gold will remain buried.
    \item What is the outcome of travelling the following route:
    \begin{itemize}
    \item $2$km on a \ealgloss{bearing}{azymut} $179\degree$
    \item $2$km on a \ealgloss{bearing}{azymut} of $000\degree$
    \item $1$km on a \ealgloss{bearing}{azymut} of $182\degree$ 
    \item $1$km on a \ealgloss{bearing}{azymut} of $000\degree$
    \end{itemize}
    \item What is the outcome of travelling the following route:
    \begin{itemize}
        \item $2$km on a \ealgloss{bearing}{azymut} $091\degree$
        \item $2$km on a \ealgloss{bearing}{azymut} of $270\degree$
        \item $1$km on a \ealgloss{bearing}{azymut} of $088\degree$
        \item  $1$km on a \ealgloss{bearing}{azymut} of $270\degree$
    \end{itemize}
    \item Use the previous two questions to show that a route leading to the gold is possible - how long will this route take?
    \item Show that it is possible to plot a route to within $1$m of any location $(x,y)$
    \item Show that it is possible to plot a route with at most $6$ legs to within $1$m of any location $(x,y)$
\end{enumerate}
\end{ealglossed}

\vspace{1em}
\section{Challenge II}

Now we restrict our attention to routes that have only two legs, thus take the following form:

\begin{ealglossed}
\begin{itemize}
    \item Travel $M$km on a \ealgloss{bearing}{azymut} of $ABC\degree$
    \item Travel $N$km on a \ealgloss{bearing}{azymut} of $XYZ\degree$
\end{itemize}
\end{ealglossed}

Where $000\degree \leq ABC, XYZ < 360\degree$, and $M,N \in \mathbb{N}$

\begin{ealglossed}
\begin{enumerate}
    \item Consider the case that $M>6000$, what restriction does this place on $XYZ\degree$ to return within $100$km of the start location?
    \item Use the \ealgloss{cosine rule}{twierdzenie cosinusów} to determine the \ealgloss{range}{rozstęp} of $N$ values possible given $M$, to ensure that the end of the journey lies between $0$ and $100$km from the start
    \item Thus establish an \ealgloss{upper bound}{ograniczenie górne} on the number of possible two leg journeys that we need to consider
    \item Does the conclusion from the previous section still necessarily hold?
    \item Use the following table to refine the \ealgloss{upper bound}{ograniczenie górne} of the number of possible two leg journeys to consider:

\begin{table}[h]
\centering
\begin{tabularx}{\textwidth}{|l|X|X|X|X|X|}
\hline
 & $ABC^{\circ}$ Option count &  $M$ in \ealgloss{range}{rozstęp} &  $XYZ^{\circ}$ Option count per $M$ value (\ealgloss{upper bound}{ograniczenie górne}) & $N$ Option count per M value (\ealgloss{upper bound}{ograniczenie górne}) & Total journey option count for this $M$ \ealgloss{range}{rozstęp} \\
\hline
$M \le 100$            & & & & & \\
\hline
$100 < M \le 500$      & & & & & \\
\hline
$500 < M \le 2000$     & & & & & \\
\hline
$2000 < M \le 6000$    & & & & & \\
\hline
\end{tabularx}
\end{table}

\item Consider the journey:
\begin{itemize}
    \item Travel $N$km on a \ealgloss{bearing}{azymut} of $XYZ\degree$
    \item Travel $M$km on a \ealgloss{bearing}{azymut} of $ABC\degree$
\end{itemize}
and thus establish and \ealgloss{upper bound}{ograniczenie górne} $U$, on the number of two leg journeys that finish within $100$km of the start point, with $U < 10^{10}$

\item Consider the \ealgloss{area}{pole} in square meters of a circle \ealgloss{radius}{promień} $100$km, thus show that there are some locations that cannot be reached within $1$m, following a two leg journey
\item Use a computer to determine if the gold can be reached via a two leg journey
\item If the gold can be reached, use a computer to find the closest point to the gold that cannot be reached; if the gold cannot be reached, find the point closest to the gold that can

\end{enumerate}
\end{ealglossed}

\end{document}